\chapter{实战：腾讯轻量云服务器部署与使用 OpenClaw}

\vspace{0.5em}
\noindent\textit{前九章我们系统学习了 OpenClaw 的十大核心概念：Gateway、Agent、Channel、Binding、Session、Plugin、Skill、Memory 和 Sandbox。
本章将全部知识付诸实践——在腾讯轻量云服务器上从零开始部署 OpenClaw，完成 API 配置、Agent 管理、
渠道接入等全流程操作。所有命令和输出均来自真实环境，读者可直接参考复现。}
\vspace{0.5em}

\begin{quote}
本章基于\textbf{腾讯轻量应用服务器}的实际操作记录，从零开始完成 OpenClaw 的部署、API 配置、Agent 管理，并通过 \texttt{openclaw tui} 交互模式展示日常使用方法。所有命令和输出均来自真实环境，读者可直接参考复现。
\end{quote}

\section{远程服务器部署}


\subsection{环境信息}

本次实战使用的是一台\textbf{腾讯轻量应用服务器}（Lighthouse），配置如下：

\begin{table}[h]
\centering
\small
\begin{tabular}{|l|p{0.7\textwidth}|}
\hline
项目 & 值 \\
\hline
云服务商 & 腾讯云 轻量应用服务器（Lighthouse） \\
\hline
服务器公网 IP & \texttt{your-server-ip} \\
\hline
用户名 & ubuntu \\
\hline
操作系统 & Ubuntu 24.04 LTS (Noble Numbat) \\
\hline
内核版本 & Linux 6.8.0-51-generic (x86\_64) \\
\hline
CPU & Intel Xeon Platinum 8255C @ 2.50GHz $\times$ 2 核 \\
\hline
内存 & 2 GB \\
\hline
磁盘 & 40 GB SSD（已用 17 GB，剩余 21 GB） \\
\hline
Node.js 版本 & v22.22.1 \\
\hline
OpenClaw 版本 & 2026.3.2 \\
\hline
私钥文件 & \texttt{\textasciitilde{}/.ssh/<your-key>.pem} \\
\hline
\end{tabular}
\end{table}

\begin{quote}
\textbf{提示}：腾讯轻量云服务器是非常适合部署 OpenClaw 的低成本方案。2核2GB 的最低配置即可流畅运行 OpenClaw Gateway 和 Agent，月费仅几十元。
\end{quote}

\subsection{SSH 连接}


从本地 Mac/Linux 终端连接到腾讯轻量云服务器：

\begin{lstlisting}[language=bash]
# 私钥权限必须为 600，否则 SSH 会拒绝连接
chmod 600 ~/.ssh/<your-key>.pem

# 连接腾讯轻量云服务器
ssh -i ~/.ssh/<your-key>.pem ubuntu@your-server-ip
\end{lstlisting}

\begin{quote}
\textbf{注意}：如果私钥权限过宽（如 644），SSH 会报错 \texttt{Permissions 0644 for key are too open}，必须先修复权限。腾讯轻量云默认使用密钥对登录，在控制台创建密钥对后下载 \texttt{.pem} 文件即可。
\end{quote}

连接成功后，可以看到 Ubuntu 24.04 LTS 的欢迎信息：

\begin{lstlisting}[language=bash]
ubuntu@VM-4-11-ubuntu:~$ uname -a
Linux VM-4-11-ubuntu 6.8.0-51-generic #52-Ubuntu SMP PREEMPT_DYNAMIC
  Thu Dec 5 13:09:44 UTC 2024 x86_64 x86_64 x86_64 GNU/Linux

ubuntu@VM-4-11-ubuntu:~$ free -h
               total        used        free
Mem:           1.9Gi       913Mi       300Mi
Swap:          1.9Gi       273Mi       1.7Gi
\end{lstlisting}

\subsection{安装 OpenClaw}

\subsubsection{系统要求}

\begin{itemize}
  \item \textbf{Node.js 22+}（安装器脚本会自动检测并安装）
  \item macOS、Linux 或通过 WSL2 的 Windows
  \item \texttt{pnpm} 仅在从源代码构建时需要
\end{itemize}

\subsubsection{方式一：安装器脚本（推荐）}

安装器脚本会自动处理 Node.js 检测、安装以及新手引导，一步到位：

\begin{lstlisting}[language=bash]
curl -fsSL https://openclaw.ai/install.sh | bash
\end{lstlisting}

安装器会自动执行以下步骤：
\begin{enumerate}
  \item 检测操作系统（macOS / Linux / WSL）
  \item 确保 Node.js 22+ 已安装（Linux 通过 NodeSource 安装）
  \item 确保 Git 已安装
  \item 通过 npm 全局安装 OpenClaw
  \item 运行新手引导（onboarding）
\end{enumerate}

如需跳过新手引导（适用于自动化部署场景）：

\begin{lstlisting}[language=bash]
curl -fsSL https://openclaw.ai/install.sh | bash -s -- --no-onboard
\end{lstlisting}

\subsubsection{方式二：全局 npm 安装（手动）}

如果服务器上已经有 Node.js 22+，可以直接通过 npm 安装：

\begin{lstlisting}[language=bash]
# 确认 Node.js 版本 >= 22
node -v

# 全局安装 OpenClaw
npm install -g openclaw@latest
\end{lstlisting}

\subsubsection{方式三：从源代码构建}

适用于需要自定义修改或参与开发的场景：

\begin{lstlisting}[language=bash]
git clone https://github.com/openclaw/openclaw.git
cd openclaw
pnpm install
pnpm build

# 将 openclaw 命令链接到全局
pnpm link --global
\end{lstlisting}

\subsubsection{安装后验证}

以下是在腾讯轻量云服务器上的实际验证结果：

\begin{lstlisting}[language=bash]
# 检查安装是否成功
ubuntu@VM-4-11-ubuntu:~$ openclaw --version
2026.3.2

# 检查 Node.js 版本
ubuntu@VM-4-11-ubuntu:~$ node -v
v22.22.1

# 验证配置文件
ubuntu@VM-4-11-ubuntu:~$ openclaw config validate
Config valid: ~/.openclaw/openclaw.json

# 检查 Gateway 健康状态
ubuntu@VM-4-11-ubuntu:~$ openclaw gateway health
Gateway Health
OK (0ms)
Feishu: ok
\end{lstlisting}

\begin{quote}
\textbf{提示}：如果安装完成后在新终端中找不到 \texttt{openclaw} 命令，通常是 Node/npm 的 PATH 配置问题。可以将 npm 全局 bin 目录加入 PATH：
\begin{lstlisting}[language=bash]
export PATH="$(npm prefix -g)/bin:$PATH"
\end{lstlisting}
建议将上述命令添加到 \texttt{\textasciitilde{}/.bashrc} 或 \texttt{\textasciitilde{}/.zshrc} 中以持久生效。
\end{quote}

\subsubsection{低配服务器优化}

腾讯轻量云的 2核2GB 属于较低配置，可以通过以下环境变量优化 OpenClaw 的启动速度：

\begin{lstlisting}[language=bash]
# 开启 Node.js 编译缓存，加速重复启动
export NODE_COMPILE_CACHE=/var/tmp/openclaw-compile-cache
mkdir -p /var/tmp/openclaw-compile-cache

# 禁止自动重启，减少启动开销
export OPENCLAW_NO_RESPAWN=1

# 写入 ~/.bashrc 持久生效
echo 'export NODE_COMPILE_CACHE=/var/tmp/openclaw-compile-cache' >> ~/.bashrc
echo 'export OPENCLAW_NO_RESPAWN=1' >> ~/.bashrc
\end{lstlisting}

\begin{quote}
\textbf{说明}：这两项优化来自 \texttt{openclaw doctor} 的建议。\texttt{NODE\_COMPILE\_CACHE} 会缓存 V8 编译结果，使后续启动加速约 30--50\%；\texttt{OPENCLAW\_NO\_RESPAWN} 避免不必要的进程重启。
\end{quote}

安装完成后，通过 \texttt{openclaw} 命令即可使用。

\hrulefill

\section{API 配置}


\subsection{配置方式}


推荐使用 \texttt{openclaw config set} 命令逐项配置，避免手动编辑 JSON 格式出错。

以下以 \textbf{DeepSeek} 为例（本服务器实际使用的 Provider）：

\begin{lstlisting}[language=bash]
# 设置 DeepSeek API Provider
openclaw config set secrets.providers.deepseek.apiKey "sk-your-deepseek-api-key"
openclaw config set secrets.providers.deepseek.apiUrl "https://api.deepseek.com/v1/"

# 设置模型（注意：模型 ID 必须是 API 实际支持的名称）
openclaw config set models.deepseek/deepseek-chat.provider "deepseek"
openclaw config set models.deepseek/deepseek-chat.api "openai-completions"

# 设置默认模型
openclaw config set defaults.model "deepseek/deepseek-chat"

# 设置网关模式为本地
openclaw config set gateway.mode "local"
\end{lstlisting}

\begin{quote}
\textbf{提示}：DeepSeek API 的模型 ID 为 \texttt{deepseek-chat}（对应 DeepSeek-V3.2），不是 \texttt{DeepSeek-V3.2}。填错会收到 \texttt{400 Model Not Exist} 错误。可通过以下命令查询实际支持的模型列表：
\begin{lstlisting}[language=bash]
curl -s https://api.deepseek.com/v1/models \
  -H "Authorization: Bearer sk-xxx" | python3 -m json.tool
\end{lstlisting}
\end{quote}

\subsection{配置文件结构}


配置文件位于 \texttt{\textasciitilde{}/.openclaw/openclaw.json}，以下是本服务器的\textbf{实际配置结构}（敏感信息已脱敏）：

\begin{lstlisting}[language=json]
{
  "meta": {
    "lastTouchedVersion": "2026.3.2",
    "lastTouchedAt": "2026-03-08T07:51:33.925Z"
  },
  "models": {
    "providers": {
      "deepseek": {
        "baseUrl": "https://api.deepseek.com/v1/",
        "apiKey": "sk-xxxx...xxxx",
        "models": [{
          "id": "deepseek-chat",
          "name": "DeepSeek-V3.2",
          "api": "openai-completions"
        }]
      }
    }
  },
  "agents": {
    "defaults": {
      "model": { "primary": "deepseek/deepseek-chat" },
      "models": { "deepseek/deepseek-chat": {} },
      "compaction": { "mode": "safeguard" }
    },
    "list": [
      { "id": "main" },
      {
        "id": "coder",
        "name": "coder",
        "workspace": "/home/ubuntu/.openclaw/workspace-coder",
        "model": "deepseek/deepseek-chat"
      }
    ]
  },
  "channels": {
    "feishu": {
      "enabled": true,
      "appId": "cli_xxxxxxxxxxxx",
      "appSecret": "***",
      "connectionMode": "websocket",
      "domain": "feishu",
      "groupPolicy": "open",
      "dmPolicy": "open"
    }
  },
  "gateway": {
    "mode": "local",
    "auth": { "mode": "token", "token": "***" }
  },
  "plugins": {
    "entries": { "feishu": { "enabled": true } },
    "installs": {
      "feishu": {
        "source": "npm",
        "spec": "@openclaw/feishu",
        "version": "2026.3.7"
      }
    }
  }
}
\end{lstlisting}

\subsection{关键字段说明}


\begin{table}[h]
\centering
\small
\begin{tabular}{|l|p{0.7\textwidth}|}
\hline
字段 & 说明 \\
\hline
\texttt{models.providers.<name>.apiKey} & API 密钥 \\
\hline
\texttt{models.providers.<name>.baseUrl} & API 端点地址（注意结尾带 \texttt{/v1/}） \\
\hline
\texttt{models.providers.<name>.models[].id} & API 实际模型 ID（必须精确，如 \texttt{deepseek-chat}） \\
\hline
\texttt{models.providers.<name>.models[].api} & API 协议类型（如 \texttt{openai-completions}） \\
\hline
\texttt{agents.defaults.model.primary} & 默认使用的模型（格式：\texttt{provider/model-id}） \\
\hline
\texttt{agents.list[]} & Agent 列表，每个 Agent 可指定独立的 workspace 和模型 \\
\hline
\texttt{channels.<channel>} & 渠道配置（如飞书、Telegram、Discord） \\
\hline
\texttt{gateway.mode} & 网关模式，\texttt{local} 为本地模式 \\
\hline
\texttt{plugins} & 插件配置（如 \texttt{@openclaw/feishu}） \\
\hline
\end{tabular}
\end{table}

\subsection{常见配置错误}


\begin{table}[h]
\centering
\small
\begin{tabular}{|l|l|l|}
\hline
错误 & 原因 & 解决方案 \\
\hline
\texttt{config validate} 失败 & api 类型写成了 \texttt{openai-chat} & 改为 \texttt{openai-completions} \\
\hline
\texttt{400 Model Not Exist} & 模型 ID 与 API 实际不匹配 & 用 \texttt{/v1/models} 接口查询实际 ID \\
\hline
\texttt{No API provider registered} & 模型配置缺少 \texttt{api} 字段 & 添加 \texttt{api: "openai-completions"} \\
\hline
网关未启动 & 未设置 \texttt{gateway.mode} 或未启动 & 设置为 \texttt{local}，执行 \texttt{openclaw gateway} \\
\hline
\texttt{duplicate plugin id} 警告 & 插件重复注册 & 检查 \texttt{plugins.entries} 配置是否重复 \\
\hline
\end{tabular}
\end{table}

\begin{quote}
\textbf{实际案例}：本服务器上即出现了 \texttt{duplicate plugin id detected} 警告，因为飞书插件同时在内置和 \texttt{extensions/feishu} 目录中被加载。该警告不影响功能使用，但会在每次执行命令时显示。
\end{quote}

\subsection{启动网关}


\begin{lstlisting}[language=bash]
# 前台运行（调试用）
openclaw gateway

# 后台运行（生产推荐）
nohup openclaw gateway > /dev/null 2>&1 &
\end{lstlisting}

\subsection{验证配置}


以下是本服务器上的实际验证过程：

\begin{lstlisting}[language=bash]
# 验证配置是否正确
ubuntu@VM-4-11-ubuntu:~$ openclaw config validate
Config valid: ~/.openclaw/openclaw.json

# 检查 Gateway 健康状态
ubuntu@VM-4-11-ubuntu:~$ openclaw gateway health
Gateway Health
OK (0ms)
Feishu: ok

# 测试 Agent 调用（发送一条消息，查看是否能正常回复）
ubuntu@VM-4-11-ubuntu:~$ openclaw agent --agent main \
  --message '你好，简单介绍一下你自己，50字以内'
我是你的AI助手OpenClaw，能帮你处理文档、搜索信息、
管理文件等任务。简洁高效，随时待命。
\end{lstlisting}

\begin{quote}
\textbf{说明}：\texttt{Gateway Health} 显示 \texttt{OK} 并且 \texttt{Feishu: ok}，说明本地网关和飞书渠道均已正常工作。Agent 也能正常回复，配置完成！
\end{quote}

\hrulefill

\section{\texttt{\textasciitilde{}/.openclaw/} 目录结构详解}


\begin{quote}
以下为服务器上 \texttt{\textasciitilde{}/.openclaw/} 的\textbf{实际完整目录树}（基于腾讯轻量云服务器实际环境采集，OpenClaw 2026.3.2）。
\end{quote}

\begin{lstlisting}
~/.openclaw/
+-- openclaw.json              # 主配置文件（API密钥、模型、网关、渠道、插件等）
+-- openclaw.json.bak          # 配置文件自动备份（最多保留5份）
+-- openclaw.json.bak.1
+-- openclaw.json.bak.2
+-- openclaw.json.bak.3
+-- openclaw.json.bak.4
+-- update-check.json          # 版本更新检查记录
|
+-- agents/                    # Agent 注册目录（每个 Agent 一个子目录）
|   +-- main/                  # 默认 Agent "main"
|   |   +-- agent/
|   |   |   +-- models.json    # 运行时模型信息快照
|   |   +-- sessions/
|   |       +-- sessions.json  # 会话索引（29个会话）
|   |       +-- <uuid>.jsonl   # 具体会话的消息历史
|   +-- coder/                 # 编程 Agent（独立工作空间）
|   +-- work/                  # 工作 Agent
|
+-- browser/                   # 浏览器自动化
|   +-- openclaw/
|
+-- canvas/                    # Canvas 画布功能
|   +-- index.html
|
+-- completions/               # Shell 自动补全脚本
|   +-- openclaw.bash
|   +-- openclaw.zsh
|   +-- openclaw.fish
|   +-- openclaw.ps1
|
+-- cron/                      # 定时任务
|   +-- jobs.json              # 定时任务注册表（如喝水提醒）
|   +-- runs/                  # 任务执行历史
|
+-- delivery-queue/            # 消息投递队列（发往渠道的消息队列）
|   +-- failed/                # 投递失败的消息
|
+-- devices/                   # 设备管理
|   +-- paired.json
|   +-- pending.json
|
+-- extensions/                # 外部插件目录
|   +-- feishu/                # 飞书插件（@openclaw/feishu v2026.3.7）
|
+-- feishu/                    # 飞书渠道运行时数据
|   +-- dedup/                 # 消息去重缓存
|
+-- identity/                  # 身份与认证
|   +-- device.json
|   +-- device-auth.json
|
+-- logs/                      # 运行日志
|   +-- config-audit.jsonl     # 配置变更审计日志
|
+-- memory/                    # 记忆存储
|   +-- main.sqlite            # 默认Agent的SQLite记忆数据库
|
+-- workspace/                 # 默认工作空间（main Agent）
|   +-- .git/
|   +-- .clawhub/lock.json
|   +-- .openclaw/workspace-state.json
|   +-- SOUL.md                # 核心人设定义
|   +-- AGENTS.md              # Agent工作流程定义
|   +-- IDENTITY.md            # Agent身份信息
|   +-- USER.md                # 用户画像（已记录：Luck, Asia/Shanghai）
|   +-- BOOTSTRAP.md           # 首次启动引导
|   +-- HEARTBEAT.md           # 心跳任务定义
|   +-- TOOLS.md               # 本地工具备忘
|   +-- memory/                # 日志记忆
|   +-- skills/
|       +-- xiaohongshu-mcp/   # 已安装：小红书自动化 Skill
|
+-- workspace-coder/           # coder Agent 独立工作空间
    +-- .git/
    +-- .openclaw/
    +-- SOUL.md                # coder 的独立人设
    +-- memory/
\end{lstlisting}

\begin{quote}
\textbf{与默认安装的差异}：本服务器安装了飞书插件（\texttt{extensions/feishu/}）、配置了多个 Agent（main/coder/work）、启用了定时任务（\texttt{cron/}）、并安装了小红书 Skill。这些是日常使用中逐步添加的，新安装的 OpenClaw 目录会更简洁。
\end{quote}

\subsection{根级文件详解}


\subsubsection{openclaw.json --- 主配置文件}


系统的核心配置，包含以下顶级字段：

\begin{table}[h]
\centering
\small
\begin{tabular}{|l|p{0.7\textwidth}|}
\hline
字段 & 说明 \\
\hline
\texttt{secrets.providers} & API Provider 的密钥和端点 \\
\hline
\texttt{models} & 可用模型列表及其 Provider 映射、API 协议类型 \\
\hline
\texttt{defaults} & 默认模型、默认 Agent 等全局缺省值 \\
\hline
\texttt{gateway} & 网关模式（\texttt{local}/\texttt{cloud}）和端口配置 \\
\hline
\texttt{channels} & 渠道连接配置（Telegram/飞书/WhatsApp等） \\
\hline
\end{tabular}
\end{table}

\subsubsection{openclaw.json.bak* --- 配置自动备份}


每次通过 \texttt{openclaw config set} 或其他方式修改配置时，系统会自动备份上一版配置。最多保留 \textbf{5 份}备份（\texttt{.bak} \textasciitilde{} \texttt{.bak.4}），按时间倒序排列，最新的为 \texttt{.bak}。

\begin{quote}
如果配置改坏了，可以直接 \texttt{cp openclaw.json.bak openclaw.json} 回退。
\end{quote}

\subsubsection{update-check.json --- 版本检查记录}


\begin{lstlisting}[language=json]
{ "lastCheckedAt": "2026-03-06T14:18:59.388Z" }
\end{lstlisting}

记录上次检查 OpenClaw 新版本的时间，避免频繁请求更新服务器。

\subsection{\texttt{agents/} --- Agent 注册与会话存储}


每个 Agent 拥有独立的子目录，内含运行时数据：

\begin{lstlisting}
agents/main/
+-- agent/
|   +-- models.json        # 运行时模型快照
+-- sessions/
    +-- sessions.json      # 会话索引
    +-- <uuid>.jsonl       # 会话消息历史
\end{lstlisting}

\subsubsection{agent/models.json --- 运行时模型信息}


Gateway 启动时从主配置解析生成，包含模型的完整信息：

\begin{lstlisting}[language=json]
{
  "providers": {
    "<provider>": {
      "baseUrl": "https://<api-endpoint>/v1/",
      "apiKey": "***",
      "models": [{
        "id": "Kimi-K2.5",
        "api": "openai-completions",
        "reasoning": false,
        "input": ["text"],
        "contextWindow": 200000,
        "maxTokens": 8192,
        "cost": { "input": 0, "output": 0, "cacheRead": 0, "cacheWrite": 0 }
      }]
    }
  }
}
\end{lstlisting}

\subsubsection{sessions/sessions.json --- 会话索引}


记录该 Agent 下所有会话的元数据，键名格式为 \texttt{agent:<agent名>:<会话名>}：

\begin{lstlisting}[language=json]
{
  "agent:main:main": {
    "sessionId": "<session-uuid>",
    "sessionFile": "~/.openclaw/agents/main/sessions/<session-uuid>.jsonl",
    "model": "Kimi-K2.5",
    "modelProvider": "<provider>",
    "contextTokens": 200000,
    "inputTokens": 32,
    "outputTokens": 162,
    "totalTokens": 13050,
    "cacheRead": 13018,
    "skillsSnapshot": { }
  }
}
\end{lstlisting}

\begin{quote}
可以从这里看到每个会话的模型、Token 用量、关联的 Skills 等信息。
\end{quote}

\subsubsection{sessions/<uuid>.jsonl --- 会话消息历史}


以 JSONL 格式（每行一条 JSON）存储完整的消息交互记录，包括用户消息、Agent 回复、工具调用等。文件名为会话的 UUID。

\subsection{\texttt{canvas/} --- Canvas 画布}


\begin{lstlisting}
canvas/
+-- index.html
\end{lstlisting}

提供一个本地 Web 可视化画布页面，Agent 可以在上面渲染图表、代码预览等富内容。通过浏览器访问。

\subsection{\texttt{completions/} --- Shell 自动补全}


\begin{lstlisting}
completions/
+-- openclaw.bash
+-- openclaw.zsh
+-- openclaw.fish
+-- openclaw.ps1
\end{lstlisting}

安装时自动生成的各 Shell 补全脚本，使 \texttt{openclaw} 命令支持 Tab 自动补全。加载方式：

\begin{lstlisting}[language=bash]
# Bash
source ~/.openclaw/completions/openclaw.bash

# Zsh
source ~/.openclaw/completions/openclaw.zsh
\end{lstlisting}

\subsection{\texttt{cron/} --- 定时任务}


\begin{lstlisting}[language=json]
// cron/jobs.json --- 本服务器实际配置
{
  "version": 1,
  "jobs": [
    {
      "id": "0cf922df-c276-4895-b2cb-6b98bc33706a",
      "agentId": "main",
      "name": "drink-water-reminder",
      "description": "每10分钟提醒Luck喝水",
      "enabled": false,
      "schedule": {
        "kind": "every",
        "everyMs": 600000
      },
      "sessionTarget": "isolated",
      "payload": {
        "kind": "agentTurn",
        "message": "执行喝水提醒脚本"
      },
      "delivery": {
        "mode": "announce",
        "channel": "last",
        "to": "ou_xxxxxxxxxxxxxxxxxxxxxxxxxxxxxxxx"
      },
      "state": {
        "lastRunStatus": "ok",
        "lastDelivered": true,
        "lastDeliveryStatus": "delivered"
      }
    }
  ]
}
\end{lstlisting}

注册 Agent 的定时任务。上述示例是本服务器上实际配置的「喝水提醒」任务：每 10 分钟触发一次，通过飞书渠道发送提醒消息。\texttt{delivery} 字段指定了消息投递目标（渠道和用户 ID）。

\begin{quote}
\textbf{提示}：定时任务通常由 Agent 在对话中自动创建（通过 \texttt{cron} Skill），也可以手动编辑 \texttt{jobs.json}。\texttt{enabled: false} 表示任务已暂停，改为 \texttt{true} 即可恢复执行。
\end{quote}

\subsection{\texttt{devices/} --- 设备管理}


\begin{lstlisting}
devices/
+-- paired.json     # 已配对设备
+-- pending.json    # 待配对设备
\end{lstlisting}

用于管理通过 WebSocket 连接到 Gateway 的外部设备（手机、平板等 Node 节点）。设备可以暴露硬件能力（摄像头、GPS、通知等），Agent 通过 Gateway 远程调用这些能力。

\subsection{\texttt{identity/} --- 身份认证}


\begin{lstlisting}
identity/
+-- device.json         # 当前设备标识信息
+-- device-auth.json    # 设备认证凭据
\end{lstlisting}

标识当前安装实例的唯一身份，用于多设备场景下的认证与区分。

\subsection{\texttt{logs/} --- 审计日志}


\begin{lstlisting}
logs/
+-- config-audit.jsonl
\end{lstlisting}

以 JSONL 格式记录每次配置变更的完整审计信息，包括：

\begin{itemize}
  \item \textbf{时间戳}（\texttt{ts}）
  \item \textbf{触发命令}（\texttt{argv}），如 \texttt{openclaw config set gateway.mode local}
  \item \textbf{变更前后的配置哈希}（\texttt{previousHash} / \texttt{nextHash}）
  \item \textbf{文件大小变化}（\texttt{previousBytes} / \texttt{nextBytes}）
  \item \textbf{可疑操作标记}（\texttt{suspicious}）
\end{itemize}

\begin{quote}
便于排查配置问题和追溯变更历史。
\end{quote}

\subsection{\texttt{memory/} --- 记忆数据库}


\begin{lstlisting}
memory/
+-- main.sqlite
\end{lstlisting}

使用 SQLite 数据库持久化 Agent 的长期记忆。文件名对应 Agent 名称（\texttt{main.sqlite} 对应 \texttt{main} Agent）。存储内容包括：

\begin{itemize}
  \item 对话摘要
  \item 用户偏好
  \item 重要事件
  \item 提取的关键信息
\end{itemize}

\subsection{\texttt{workspace/} --- 工作空间（人设与知识库）}


这是 Agent 的"灵魂所在"，包含人设定义、工作流程和已安装的 Skills：

\begin{lstlisting}
workspace/
+-- .git/                          # 自带Git版本控制
+-- .clawhub/                      # ClawHub 包管理
|   +-- lock.json                  # Skills 版本锁定文件
+-- .openclaw/workspace-state.json # 工作空间状态
+-- SOUL.md          # 核心人设
+-- AGENTS.md        # 工作流程
+-- IDENTITY.md      # 身份信息
+-- USER.md          # 用户画像
+-- BOOTSTRAP.md     # 首次引导
+-- HEARTBEAT.md     # 心跳任务
+-- TOOLS.md         # 本地工具备忘
|
+-- skills/          # 通过 ClawHub 安装的工作空间级 Skills
    +-- <skill-name>/
        +-- .clawhub/origin.json   # 安装来源信息
        +-- _meta.json             # Skill 元数据
        +-- SKILL.md               # Skill 定义文件
        +-- scripts/               # Skill 附带脚本（可选）
\end{lstlisting}

\subsubsection{SOUL.md --- 核心人设定义}


定义 Agent 的性格和行为准则，\textbf{是最重要的人设文件}。默认内容要点：

\begin{itemize}
  \item \textbf{真诚帮助}：跳过客套话（如"Great question!"），直接解决问题
  \item \textbf{有观点}：允许有偏好和意见，不做没有个性的搜索引擎
  \item \textbf{先自己查}：先读文件、查上下文、搜索，实在找不到再问用户
  \item \textbf{尊重隐私}：用户给了访问权限，不要滥用，隐私信息不外泄
  \item \textbf{对外谨慎，对内大胆}：对外操作（发邮件、发推）要确认，对内操作（读文件、整理）可以主动
\end{itemize}

\subsubsection{AGENTS.md --- Agent 工作流程定义}


规定 Agent 每次启动时的标准流程：

\begin{enumerate}
  \item 读取 \texttt{SOUL.md} --- 知道自己是谁
  \item 读取 \texttt{USER.md} --- 知道在帮谁
  \item 读取 \texttt{memory/YYYY-MM-DD.md}（今天+昨天）--- 获取近期上下文
  \item 如果是主会话，还要读取 \texttt{MEMORY.md} --- 长期记忆
\end{enumerate}

还定义了记忆管理规范：日志写到 \texttt{memory/YYYY-MM-DD.md}，精华提炼到 \texttt{MEMORY.md}。

\subsubsection{IDENTITY.md --- Agent 身份信息}


存储 Agent 的具体身份标识：名字、物种/形态、沟通风格、代表 emoji 等。在首次引导（Bootstrap）时由用户和 Agent 共同确定。

\subsubsection{USER.md --- 用户画像}


记录用户的基本信息和偏好。以下是本服务器上的实际内容：

\begin{lstlisting}
# USER.md - About Your Human

- **Name:** Luck
- **What to call them:** Luck
- **Timezone:** Asia/Shanghai (中国陕西)
- **Notes:** 用户来自中国陕西地区

## Context
_(用户关心什么？在做什么项目？...)_

---
The more you know, the better you can help.
But remember - you're learning about a person,
not building a dossier. Respect the difference.
\end{lstlisting}

Agent 通过持续交互不断更新此文件，以便更好地理解和服务用户。

\subsubsection{BOOTSTRAP.md --- 首次启动引导}


新安装后的"出生引导"脚本，引导 Agent 完成：

\begin{enumerate}
  \item 了解用户（名字、称呼、时区）
  \item 确定自己的身份（名字、风格、emoji）
  \item 更新 \texttt{IDENTITY.md}、\texttt{USER.md}、\texttt{SOUL.md}
  \item 可选配置消息渠道（WhatsApp/Telegram等）
  \item \textbf{完成后删除此文件}（不再需要引导）
\end{enumerate}

\subsubsection{HEARTBEAT.md --- 心跳任务}


定义 Agent 需要周期性执行的检查任务。文件为空或仅有注释时，跳过心跳调用。添加任务后，Agent 会按周期自动执行。

\subsubsection{TOOLS.md --- 本地工具备忘}


记录当前环境特有的工具和设备信息，如：SSH 主机别名、摄像头名称、TTS 语音偏好、智能设备昵称等。Skills 是通用的，TOOLS.md 是"你的环境独有的"。

\subsubsection{.clawhub/lock.json --- Skills 版本锁定}


类似 \texttt{package-lock.json}，记录通过 ClawHub 安装的所有 Skills 的版本和安装时间：

\begin{lstlisting}[language=json]
{
  "version": 1,
  "skills": {
    "xiaohongshu-mcp": {
      "version": "1.0.0",
      "installedAt": 1772864060491
    }
  }
}
\end{lstlisting}

\subsubsection{skills/ --- 已安装的工作空间级 Skills}


通过 ClawHub 市场安装的 Skills 会存放在此目录下。每个 Skill 是一个独立子目录，内部结构：

\begin{lstlisting}
skills/<skill-name>/
+-- .clawhub/
|   +-- origin.json        # 安装来源（registry、slug、版本、安装时间）
+-- _meta.json             # Skill 元数据（作者ID、slug、版本、发布时间）
+-- SKILL.md               # Skill 定义文件（名称、描述、使用指南）
+-- scripts/               # Skill 附带的可执行脚本（可选）
\end{lstlisting}

\textbf{origin.json} --- 记录 Skill 从哪里安装的：

\begin{lstlisting}[language=json]
{
  "version": 1,
  "registry": "https://clawhub.ai",
  "slug": "xiaohongshu-mcp",
  "installedVersion": "1.0.0",
  "installedAt": 1772864060489
}
\end{lstlisting}

\textbf{\_meta.json} --- Skill 的发布元数据：

\begin{lstlisting}[language=json]
{
  "ownerId": "<owner-id>",
  "slug": "xiaohongshu-mcp",
  "version": "1.0.0",
  "publishedAt": 1769938539127
}
\end{lstlisting}

\textbf{SKILL.md} --- Skill 的核心定义文件，包含 YAML frontmatter（名称、描述、触发条件）和 Markdown 正文（使用指南、命令说明）。

\begin{quote}
\textbf{与内置 Skills 的区别}：内置 Skills 在 \texttt{/usr/lib/node\_modules/openclaw/skills/} 下，不含 \texttt{.clawhub/} 和 \texttt{\_meta.json}；通过 ClawHub 安装的 Skills 带有完整的来源追踪信息。
\end{quote}

\subsubsection{实例：已安装的 xiaohongshu-mcp Skill}


当前服务器上已安装了小红书自动化 Skill，它提供：

\begin{table}[h]
\centering
\small
\begin{tabular}{|l|l|l|}
\hline
功能 & 命令 & 说明 \\
\hline
检查登录状态 & \texttt{python scripts/xhs\_client.py status} & 确认 MCP 服务器是否运行 \\
\hline
搜索笔记 & \texttt{python scripts/xhs\_client.py search "关键词"} & 按关键词搜索小红书笔记 \\
\hline
获取笔记详情 & \texttt{python scripts/xhs\_client.py detail <feed\_id> <xsec\_token>} & 获取指定笔记的完整内容和评论 \\
\hline
获取推荐流 & \texttt{python scripts/xhs\_client.py feeds} & 获取推荐信息流 \\
\hline
发布笔记 & \texttt{python scripts/xhs\_client.py publish "标题" "内容" "图片URL"} & 发布图文笔记 \\
\hline
\end{tabular}
\end{table}

\begin{quote}
使用前需先在本地启动 \texttt{xiaohongshu-mcp} 服务器（监听 \texttt{localhost:18060}）并完成扫码登录。
\end{quote}

\subsubsection{.openclaw/workspace-state.json --- 工作空间状态}


\begin{lstlisting}[language=json]
{ "version": 1, "bootstrapSeededAt": "2026-03-06T14:21:02.015Z" }
\end{lstlisting}

记录工作空间的初始化时间等状态信息。

\subsubsection{.git/ --- 版本控制}


workspace 目录自带 Git 仓库，Agent 对人设文件的修改可以通过 Git 追踪和回滚。

\hrulefill

\section{Skills 系统详解}


\subsection{Skills 三层存放路径}


Skills 分三层存放，\textbf{优先级从高到低}：

\begin{table}[h]
\centering
\small
\begin{tabular}{|l|l|l|}
\hline
层级 & 路径 & 说明 \\
\hline
\textbf{工作空间 Skills} & \texttt{\textasciitilde{}/.openclaw/workspace/skills/} & 仅当前工作空间 Agent 可用，优先级\textbf{最高} \\
\hline
\textbf{全局共享 Skills} & \texttt{\textasciitilde{}/.openclaw/skills/} & 所有 Agent 共享的自定义 Skills \\
\hline
\textbf{内置 Skills} & \texttt{/usr/lib/node\_modules/openclaw/skills/} & 随 OpenClaw 安装自带，优先级\textbf{最低} \\
\hline
\end{tabular}
\end{table}

\begin{quote}
同名 Skill 时，高优先级路径的会覆盖低优先级的。
\end{quote}

\subsection{Skill 的目录结构}


每个 Skill 就是一个目录，核心是 \texttt{SKILL.md} 文件。根据来源不同，目录结构有所差异：

\textbf{手动创建 / 内置 Skill（最简结构）}：

\begin{lstlisting}
weather/
+-- SKILL.md       # Skill 定义文件
\end{lstlisting}

\textbf{通过 ClawHub 市场安装的 Skill（完整结构）}：

\begin{lstlisting}
xiaohongshu-mcp/
+-- .clawhub/
|   +-- origin.json        # 安装来源（registry URL、版本、安装时间）
+-- _meta.json             # 发布元数据（作者ID、slug、版本、发布时间）
+-- SKILL.md               # Skill 定义文件
+-- scripts/               # 附带的可执行脚本（可选）
    +-- xhs_client.py
\end{lstlisting}

\texttt{SKILL.md} 由两部分组成：

\begin{itemize}
  \item \textbf{Frontmatter}（YAML 头部）：声明名称、描述、依赖要求等元数据
  \item \textbf{正文}（Markdown）：详细的使用指南和命令说明
\end{itemize}

示例（\texttt{weather} Skill）：

\begin{lstlisting}[language=yaml]
---
name: weather
description: "Get current weather and forecasts via wttr.in or Open-Meteo..."
homepage: https://wttr.in/:help
metadata: { "openclaw": { "emoji": "...", "requires": { "bins": ["curl"] } } }
---

# Weather Skill
## Commands
...
\end{lstlisting}

\subsection{Skill 的依赖检测与激活}


OpenClaw 启动时会自动检测每个 Skill 的依赖是否满足，\textbf{缺依赖的 Skill 不会被加载}。

依赖类型有四种：

\begin{table}[h]
\centering
\small
\begin{tabular}{|l|l|l|}
\hline
依赖类型 & 说明 & 示例 \\
\hline
\texttt{bins} & 需要系统中安装的命令行工具 & \texttt{curl}、\texttt{gh}、\texttt{rg}、\texttt{clawhub} \\
\hline
\texttt{env} & 需要设置的环境变量 & \texttt{OPENAI\_API\_KEY}、\texttt{GEMINI\_API\_KEY} \\
\hline
\texttt{os} & 要求特定操作系统 & \texttt{darwin}（仅 macOS） \\
\hline
\texttt{config} & 需要在配置文件中设置的字段 & \texttt{channels.discord.token} \\
\hline
\end{tabular}
\end{table}

\textbf{激活方法就是满足对应依赖}，例如：

\begin{lstlisting}[language=bash]
# clawhub skill 缺 clawhub 命令 -> 安装即可激活
sudo npm install -g clawhub

# github/gh-issues skill 缺 gh 命令
sudo apt install gh

# session-logs skill 缺 rg (ripgrep) 命令
sudo apt install ripgrep

# notion skill 缺环境变量
export NOTION_API_KEY="your-key"

# apple-notes skill 要求 os: darwin -> Linux 上无法激活，仅 macOS 可用
\end{lstlisting}

\subsection{Skills 管理命令}


\begin{lstlisting}[language=bash]
# 查看所有 Skills 的状态（哪些可用、哪些缺依赖）
openclaw skills check

# 列出所有 Skills
openclaw skills list

# 仅列出已就绪的 Skills
openclaw skills list --eligible

# 查看某个 Skill 的详细信息（含依赖要求和安装提示）
openclaw skills info <name>
# 示例：
openclaw skills info clawhub
# 输出：  clawhub  Missing requirements
#       Requirements:
#         Binaries:  clawhub
#       Install options:
#         -> Install clawhub (npm)
\end{lstlisting}

\subsection{安装第三方 Skills}


\subsubsection{方式一：通过 ClawHub 市场安装（推荐）}


ClawHub 是 OpenClaw 的 Skills 市场（https://clawhub.ai），类似 npm 之于 Node.js：

\begin{lstlisting}[language=bash]
# 先安装 clawhub CLI
sudo npm install -g clawhub

# 搜索、安装、同步 Skills
npx clawhub
\end{lstlisting}

安装后的文件变化：

% [Auto-converted from ASCII art to TikZ]
\begin{figure}[htbp]
\centering
\begin{tikzpicture}
  \node[dirtree] {
    workspace/\\
    +-- .clawhub/\\
    |   +-- lock.json              \# <- 新增/更新：记录已安装 Skill 版本\\
    +-- skills/\\
    +-- <skill-name>/          \# <- 新增：Skill 目录\\
    +-- .clawhub/origin.json   \# 安装来源追踪\\
    +-- \_meta.json             \# 作者、版本元数据\\
    +-- SKILL.md               \# Skill 定义\\
    +-- scripts/               \# 附带脚本（如有）\\
  };
\end{tikzpicture}
\end{figure}

\begin{quote}
通过 ClawHub 安装的 Skills 位于 \texttt{workspace/skills/} 下（工作空间级，优先级最高），并通过 \texttt{.clawhub/lock.json} 锁定版本。
\end{quote}

\subsubsection{方式二：通过 plugins 命令安装}


\begin{lstlisting}[language=bash]
# 从本地路径安装（支持 .ts/.js/.zip/.tgz）
openclaw plugins install ./my-skill.zip

# 从 npm 安装
openclaw plugins install some-openclaw-plugin

# 以符号链接方式安装（开发调试用）
openclaw plugins install ./my-skill --link

# 锁定精确版本
openclaw plugins install some-plugin --pin
\end{lstlisting}

插件管理完整命令：

\begin{table}[h]
\centering
\small
\begin{tabular}{|l|p{0.7\textwidth}|}
\hline
命令 & 说明 \\
\hline
\texttt{openclaw plugins install <path-or-spec>} & 安装插件（本地路径/npm包） \\
\hline
\texttt{openclaw plugins uninstall} & 卸载插件 \\
\hline
\texttt{openclaw plugins update} & 更新已安装的 npm 插件 \\
\hline
\texttt{openclaw plugins list} & 列出已安装插件 \\
\hline
\texttt{openclaw plugins info} & 查看插件详情 \\
\hline
\texttt{openclaw plugins enable} / \texttt{disable} & 启用/禁用插件 \\
\hline
\texttt{openclaw plugins doctor} & 排查插件加载问题 \\
\hline
\end{tabular}
\end{table}

\subsubsection{方式三：手动放置 Skill 目录}


直接创建目录并放入 \texttt{SKILL.md}：

\begin{lstlisting}[language=bash]
# 全局共享（所有 Agent 可用）
mkdir -p ~/.openclaw/skills/my-skill
# 编写 SKILL.md ...

# 工作空间专属（仅当前工作空间 Agent 可用，优先级最高）
mkdir -p ~/.openclaw/workspace/skills/my-skill
# 编写 SKILL.md ...
\end{lstlisting}

\subsection{内置 Skills 一览}


本服务器上共有 \textbf{13 个已就绪的 Skills}（包含飞书插件提供的 Skills），以及多个未满足依赖的 Skills。

\textbf{已就绪（13个）}：

\begin{table}[h]
\centering
\small
\begin{tabular}{|l|l|p{0.5\textwidth}|}
\hline
Skill & 来源 & 功能 \\
\hline
\texttt{feishu-doc} & openclaw-extra & 飞书文档读写操作 \\
\hline
\texttt{feishu-drive} & openclaw-extra & 飞书云空间文件管理 \\
\hline
\texttt{feishu-perm} & openclaw-extra & 飞书权限管理 \\
\hline
\texttt{feishu-wiki} & openclaw-extra & 飞书知识库导航 \\
\hline
\texttt{clawhub} & openclaw-bundled & ClawHub Skill 市场管理 \\
\hline
\texttt{healthcheck} & 内置 & 主机安全检查与加固 \\
\hline
\texttt{skill-creator} & 内置 & 创建或更新 Skills \\
\hline
\texttt{tmux} & 内置 & 远程控制 tmux 会话 \\
\hline
\texttt{video-frames} & 内置 & 从视频提取帧或片段 \\
\hline
\texttt{weather} & 内置 & 天气查询与预报 \\
\hline
\end{tabular}
\end{table}
\textbf{未就绪（46个，按类别）}：

\begin{table}[h]
\centering
\small
\begin{tabular}{|l|l|l|}
\hline
类别 & Skills & 典型缺失依赖 \\
\hline
AI/模型 & \texttt{gemini}, \texttt{openai-image-gen}, \texttt{openai-whisper}, \texttt{openai-whisper-api}, \texttt{oracle} & env: API Key \\
\hline
笔记/文档 & \texttt{apple-notes}, \texttt{bear-notes}, \texttt{notion}, \texttt{obsidian}, \texttt{nano-pdf}, \texttt{summarize} & bins / os: darwin \\
\hline
通讯/社交 & \texttt{discord}, \texttt{slack}, \texttt{imsg}, \texttt{bluebubbles}, \texttt{wacli}, \texttt{voice-call} & config / os: darwin \\
\hline
开发工具 & \texttt{coding-agent}, \texttt{github}, \texttt{gh-issues}, \texttt{trello}, \texttt{clawhub}, \texttt{mcporter} & bins: gh, clawhub 等 \\
\hline
多媒体 & \texttt{canvas}, \texttt{camsnap}, \texttt{peekaboo}, \texttt{songsee}, \texttt{spotify-player}, \texttt{gifgrep} & bins / os: darwin \\
\hline
智能家居 & \texttt{openhue}, \texttt{sonoscli}, \texttt{blucli} & bins \\
\hline
系统/运维 & \texttt{session-logs}, \texttt{model-usage} & bins: rg / os: darwin \\
\hline
信息获取 & \texttt{xurl}, \texttt{blogwatcher}, \texttt{gog}, \texttt{goplaces}, \texttt{sag} & bins / env \\
\hline
其他 & \texttt{1password}, \texttt{apple-reminders}, \texttt{things-mac}, \texttt{himalaya}, 等 & bins / env / os \\
\hline
\end{tabular}
\end{table}

\begin{quote}
\textbf{提示}：运行 \texttt{openclaw skills check} 可以实时查看每个 Skill 具体缺什么。\texttt{os: darwin} 标记的 Skill 仅限 macOS，在 Linux 服务器上无法激活。
\end{quote}

\hrulefill

\section{Agent 与会话层级关系}


\subsection{层级结构}


% [Auto-converted from ASCII art to TikZ]
\begin{figure}[htbp]
\centering
\begin{tikzpicture}
  \node[dirtree] {
    OpenClaw 实例\\
    +-- Agent（智能体）\\
    +-- Session 1（会话1）\\
    +-- Session 2（会话2）\\
    +-- Session N（会话N）\\
  };
\end{tikzpicture}
\end{figure}

\textbf{核心概念}：

\begin{itemize}
  \item \textbf{Agent} 是执行单元，每个 Agent 有独立的配置、人设和能力
  \item \textbf{Session} 是对话上下文，每个 Agent 下可以有多个 Session
  \item 层级关系：\texttt{Agent} $\rightarrow$ \texttt{Session}，一对多
\end{itemize}

\subsection{Agent 管理}


本服务器配置了多个 Agent，每个有不同的职责：

\begin{table}[h]
\centering
\small
\begin{tabular}{|l|l|l|}
\hline
Agent ID & 工作空间 & 用途 \\
\hline
\texttt{main} & \texttt{\textasciitilde{}/.openclaw/workspace/} & 默认主 Agent，日常对话、任务处理 \\
\hline
\texttt{coder} & \texttt{\textasciitilde{}/.openclaw/workspace-coder/} & 编程 Agent，独立工作空间和人设 \\
\hline
\texttt{work} & \texttt{\textasciitilde{}/.openclaw/workspace/} & 工作 Agent，共享默认工作空间 \\
\hline
\end{tabular}
\end{table}

\begin{lstlisting}[language=bash]
# 使用默认 Agent（main）
ubuntu@VM-4-11-ubuntu:~$ openclaw agent --agent main \
  --message '你好'
我是你的AI助手OpenClaw，能帮你处理文档、搜索信息、
管理文件等任务。简洁高效，随时待命。

# 使用 coder Agent（拥有独立的工作空间和人设）
ubuntu@VM-4-11-ubuntu:~$ openclaw agent --agent coder \
  --message '帮我写一个 Python 脚本'

# 不指定 --agent 时，默认使用 main Agent
\end{lstlisting}

\begin{quote}
\textbf{提示}：\texttt{coder} Agent 拥有独立的 \texttt{workspace-coder/} 目录，包含自己的 \texttt{SOUL.md}。这意味着它可以有与 \texttt{main} Agent 完全不同的人设和行为模式。
\end{quote}

\subsection{Session（会话）管理}


\subsubsection{开启新会话}


\textbf{方式一：通过 TUI 交互模式（推荐）}

\begin{lstlisting}[language=bash]
# 进入默认会话（main）的 TUI 交互模式
openclaw tui

# 开启一个新的命名会话（指定新名称即自动创建）
openclaw tui --session my-new-session

# 带初始消息开启新会话
openclaw tui --session work --message "帮我整理一下今天的任务"
\end{lstlisting}

\begin{quote}
\texttt{--session <key>} 是创建新会话的关键参数。不指定时默认使用 \texttt{main} 会话，指定一个\textbf{新名称}就会自动创建新会话，彼此上下文完全隔离。
\end{quote}

\textbf{方式二：通过 \texttt{agent} 命令（单次对话，非交互）}

\begin{lstlisting}[language=bash]
# 在默认会话中发一条消息（复用已有 main 会话上下文）
openclaw agent --agent main --message "你好"

# 使用指定 session-id 复用某个会话
openclaw agent --agent main --session-id <uuid> --message "继续上次的话题"
\end{lstlisting}

\begin{quote}
\textbf{注意}：\texttt{agent --session-id} 是用来\textbf{复用已有会话 UUID} 的，不是创建新命名会话。要创建全新会话，推荐使用 \texttt{openclaw tui --session <新名称>}。
\end{quote}

\subsubsection{查看会话列表}


以下是本服务器上的实际会话列表（截取部分）：

\begin{lstlisting}[language=bash]
ubuntu@VM-4-11-ubuntu:~$ openclaw sessions
Session store: ~/.openclaw/agents/main/sessions/sessions.json
Sessions listed: 29
Kind   Key                        Age    Model          Tokens (ctx %)
direct agent:main:main            1h ago deepseek-chat  19k/200k (10%)
direct agent:main:luck1           3d ago deepseek-chat  29k/200k (14%)
direct agent:main:cron:...33706a  3d ago deepseek-chat  17k/200k (9%)
  ...(省略多个 cron 任务会话)...
\end{lstlisting}

\begin{quote}
\textbf{解读}：可以看到服务器上共有 29 个会话，包括：
\begin{itemize}
  \item \texttt{agent:main:main} --- 默认主会话，使用 deepseek-chat 模型，已用 19k/200k Token
  \item \texttt{agent:main:luck1} --- 用户自定义的命名会话
  \item \texttt{agent:main:cron:*} --- 定时任务自动创建的隔离会话
\end{itemize}
\end{quote}

常用查询命令：

\begin{lstlisting}[language=bash]
# 列出所有会话
openclaw sessions

# 仅查看最近 N 分钟活跃的会话
openclaw sessions --active 120

# JSON 格式输出
openclaw sessions --json

# 查看指定 Agent 的会话
openclaw sessions --agent work

# 查看所有 Agent 的会话
openclaw sessions --all-agents
\end{lstlisting}

\subsubsection{清理会话}


\begin{lstlisting}[language=bash]
# 预览要清理的会话（不实际执行）
openclaw sessions cleanup --dry-run

# 真正清理过期会话
openclaw sessions cleanup --enforce

# 手动清空所有会话（完全重置）
echo '{}' > ~/.openclaw/agents/main/sessions/sessions.json
\end{lstlisting}

\subsubsection{TUI 完整参数}


\begin{table}[h]
\centering
\small
\begin{tabular}{|l|p{0.7\textwidth}|}
\hline
参数 & 说明 \\
\hline
\texttt{--session <key>} & 会话名称，指定新名称自动创建（默认 \texttt{main}） \\
\hline
\texttt{--message <text>} & 连接后自动发送的初始消息 \\
\hline
\texttt{--thinking <level>} & 思考等级覆盖 \\
\hline
\texttt{--deliver} & 将回复投递到渠道（默认 false） \\
\hline
\texttt{--history-limit <n>} & 加载的历史条目数（默认 200） \\
\hline
\texttt{--url <url>} & Gateway WebSocket URL \\
\hline
\texttt{--token <token>} & Gateway token \\
\hline
\end{tabular}
\end{table}

\subsection{Agent 独立人设配置}


每个 Agent \textbf{可以}拥有独立的人设，通过 \texttt{--workspace} 参数指定独立工作空间：

\begin{lstlisting}[language=bash]
# 为不同 Agent 指定不同的工作空间
openclaw agent --agent work --workspace ~/openclaw-work --message "处理工作邮件"
openclaw agent --agent life --workspace ~/openclaw-life --message "今天天气怎么样"
\end{lstlisting}

不同工作空间下可以有各自独立的 \texttt{SOUL.md}，从而实现不同 Agent 拥有不同的人设定义。

\begin{quote}
\textbf{说明}：默认情况下 \texttt{\textasciitilde{}/.openclaw/workspace/} 是共享的，所有 Agent 使用同一套人设。如果需要独立人设，需通过 \texttt{--workspace} 参数为不同 Agent 指定不同的工作空间目录。
\end{quote}

\hrulefill

\section{模型切换实战}


\subsection{切换步骤}


以从 DeepSeek-V3.2 切换到其他模型（如 Kimi-K2.5）为例：

\textbf{第一步：备份当前配置}

\begin{lstlisting}[language=bash]
cp ~/.openclaw/openclaw.json ~/.openclaw/openclaw.json.bak.old-provider
\end{lstlisting}

\textbf{第二步：修改配置文件}

可通过 \texttt{openclaw config set} 命令修改（推荐），也可通过 Python 脚本批量修改（避免手动编辑 JSON 出错）：

\begin{lstlisting}[language=bash]
# 方式一：用 openclaw config set 命令（推荐）
openclaw config set models.providers.kimi.baseUrl "https://api.moonshot.cn/v1/"
openclaw config set models.providers.kimi.apiKey "sk-your-kimi-key"
openclaw config set agents.defaults.model.primary "kimi/moonshot-v1-128k"
\end{lstlisting}

\begin{lstlisting}[language=python]
# 方式二：Python 脚本批量修改
import json, os

config_path = os.path.expanduser('~/.openclaw/openclaw.json')
c = json.load(open(config_path))

# 替换 models.providers
c['models'] = {
    'providers': {
        'kimi': {
            'baseUrl': 'https://api.moonshot.cn/v1/',
            'apiKey': 'sk-xxxxxxxxxxxxxxxx',
            'models': [{
                'id': 'moonshot-v1-128k',       # <- API 实际模型 ID
                'name': 'Kimi-K2.5',            # <- 显示名称（自定义）
                'api': 'openai-completions'
            }]
        }
    }
}

# 更新默认模型引用
c['agents']['defaults']['model'] = {'primary': 'kimi/moonshot-v1-128k'}
c['agents']['defaults']['models'] = {'kimi/moonshot-v1-128k': {}}

json.dump(c, open(config_path, 'w'), indent=2)
\end{lstlisting}

\textbf{第三步：清理运行时缓存和旧会话}

\begin{quote}
这一步至关重要！不清理会导致旧模型信息残留。
\end{quote}

\begin{lstlisting}[language=bash]
# 清理运行时模型快照（Gateway 重启后会重新生成）
rm ~/.openclaw/agents/main/agent/models.json

# 清空旧会话（旧会话上下文中包含旧模型信息，会导致新模型"自以为是"旧模型）
echo '{}' > ~/.openclaw/agents/main/sessions/sessions.json
\end{lstlisting}

\textbf{第四步：重启 Gateway}

\begin{lstlisting}[language=bash]
# 杀掉旧进程
pkill -f openclaw-gateway

# 重新启动
nohup openclaw gateway run --force > /dev/null 2>&1 &

# 验证
openclaw gateway health
openclaw status
\end{lstlisting}

\textbf{第五步：测试新模型}

\begin{lstlisting}[language=bash]
openclaw agent --agent main --message "你是什么模型?" --json
\end{lstlisting}

\begin{quote}
\textbf{实际案例}：本服务器上就是从最初的 Kimi 模型切换到 DeepSeek-V3.2 的。切换后测试：
\begin{lstlisting}[language=bash]
ubuntu@VM-4-11-ubuntu:~$ openclaw agent --agent main \
  --message '你是什么模型?'
我是 DeepSeek-V3.2，由 DeepSeek 团队开发...
\end{lstlisting}
如果 Agent 回复的仍是旧模型名称，说明没有完成下述“三处清理”。
\end{quote}

\subsection{关键注意事项}


\subsubsection{模型 ID 必须用 API 实际支持的名称}


不同 API 提供商的模型 ID 可能与官方宣传名称不同：

\begin{lstlisting}[language=bash]
# 查询 DeepSeek API 实际支持的模型列表
curl -s https://api.deepseek.com/v1/models \
  -H "Authorization: Bearer sk-xxx" | python3 -c \
  "import sys,json; [print(m['id']) for m in \
  json.load(sys.stdin).get('data',[])]"
# 输出：deepseek-chat  deepseek-reasoner

# 查询 Kimi API 实际支持的模型列表
curl -s https://api.moonshot.cn/v1/models \
  -H "Authorization: Bearer sk-xxx" | python3 -c \
  "import sys,json; [print(m['id']) for m in \
  json.load(sys.stdin).get('data',[])]"
\end{lstlisting}

\textbf{DeepSeek 实际模型 ID}：

\begin{table}[h]
\centering
\small
\begin{tabular}{|l|l|l|}
\hline
官方名称 & API 模型 ID & 说明 \\
\hline
DeepSeek-V3.2 & \texttt{deepseek-chat} & 通用对话模型 \\
\hline
DeepSeek-R1 & \texttt{deepseek-reasoner} & 推理模型 \\
\hline
\end{tabular}
\end{table}

\begin{quote}
填错模型 ID 会得到 \texttt{400 Model Not Exist} 错误。
\end{quote}

\subsubsection{配置文件中的 id 与 name 的区别}


\begin{lstlisting}[language=json]
{
  "id": "deepseek-chat",      // <- 实际发送给 API 的模型 ID（必须准确）
  "name": "DeepSeek-V3.2",    // <- 仅用于显示的友好名称（可自定义）
  "api": "openai-completions"
}
\end{lstlisting}

\subsubsection{切换模型后必须清理的三个地方}


\begin{table}[h]
\centering
\small
\begin{tabular}{|l|l|l|}
\hline
清理项 & 路径 & 原因 \\
\hline
运行时模型快照 & \texttt{agents/main/agent/models.json} & Gateway 重启后会合并新旧配置，导致旧 provider 残留 \\
\hline
旧会话 & \texttt{agents/main/sessions/sessions.json} & 旧会话上下文中包含"我是XX模型"的历史消息 \\
\hline
Gateway 进程 & \texttt{pkill -f openclaw-gateway} & 内存中的配置需要重启才能刷新 \\
\hline
\end{tabular}
\end{table}

\begin{quote}
不清理这三处，即使配置改了，Agent 还是会根据旧的上下文信息回答说自己是旧模型。
\end{quote}

\subsection{快速回退}


\begin{lstlisting}[language=bash]
# 切回之前的模型配置
cp ~/.openclaw/openclaw.json.bak.old-provider ~/.openclaw/openclaw.json
rm ~/.openclaw/agents/main/agent/models.json
echo '{}' > ~/.openclaw/agents/main/sessions/sessions.json
pkill -f openclaw-gateway
nohup openclaw gateway run --force > /dev/null 2>&1 &
\end{lstlisting}

\hrulefill

\section{实战踩坑总结}


\begin{table}[h]
\centering
\small
\begin{tabular}{|l|l|l|}
\hline
问题 & 现象 & 解决 \\
\hline
SSH 私钥权限 & \texttt{Permissions too open} 拒绝连接 & \texttt{chmod 600 \textasciitilde{}/.ssh/<your-key>.pem} \\
\hline
JSON 配置格式 & 手动编辑容易出错 & 用 \texttt{openclaw config set} 命令配置 \\
\hline
API 类型错误 & \texttt{openai-chat} 不被支持 & 改为 \texttt{openai-completions} \\
\hline
模型缺少 api 字段 & \texttt{No API provider for api: undefined} & 配置中添加 \texttt{api} 字段 \\
\hline
网关未启动 & Agent 无法调用 & \texttt{nohup openclaw gateway \&} \\
\hline
模型 ID 写错 & \texttt{400 Model Not Exist} & 用 \texttt{/v1/models} 接口查询实际 ID \\
\hline
Agent 自称旧模型 & 配置已改但回复说“我是 Kimi” & 清理三处：models.json + sessions + 重启 \\
\hline
\texttt{duplicate plugin id} & 飞书插件重复加载警告 & 检查 extensions/ 和内置插件是否重复 \\
\hline
\texttt{--session-id} 未创建新会话 & 期望新建会话但复用了旧的 & 用 \texttt{openclaw tui --session <新名称>} \\
\hline
低内存 OOM & 2GB 内存偶尔被杀 & 启用 Swap + \texttt{OPENCLAW\_NO\_RESPAWN} \\
\hline
Gateway SIGHUP 崩溃 & 发 HUP 信号后 Gateway 挂掉 & 不要用 \texttt{kill -HUP}，直接 pkill 后重启 \\
\hline
\end{tabular}
\end{table}

\clearpage

\section{附录：飞书渠道配置实战}

\subsection{配置概览}

本服务器已配置并开启了\textbf{飞书渠道}，让 OpenClaw Agent 能够通过飞书与用户交互。配置如下：

\begin{lstlisting}[language=bash]
# 1. 安装飞书插件
openclaw plugins install @openclaw/feishu

# 2. 配置飞书应用凭据
openclaw config set channels.feishu.enabled true
openclaw config set channels.feishu.appId "cli_xxxxxxxxxxxx"
openclaw config set channels.feishu.appSecret "your-app-secret"
openclaw config set channels.feishu.connectionMode "websocket"
openclaw config set channels.feishu.domain "feishu"
openclaw config set channels.feishu.groupPolicy "open"
openclaw config set channels.feishu.dmPolicy "open"
\end{lstlisting}

\subsection{验证飞书状态}

\begin{lstlisting}[language=bash]
ubuntu@VM-4-11-ubuntu:~$ openclaw gateway health
Gateway Health
OK (0ms)
Feishu: ok     # <- 飞书连接正常
\end{lstlisting}

配置完成后，可以在飞书中直接与 Agent 对话，消息会通过 WebSocket 实时传递。定时任务（如喝水提醒）也会通过飞书渠道推送消息。

\begin{quote}
\textbf{提示}：飞书渠道需要在飞书开放平台创建应用，获取 \texttt{appId} 和 \texttt{appSecret}。\texttt{groupPolicy: "open"} 表示允许所有群聊触发 Agent，\texttt{dmPolicy: "open"} 表示允许所有私聊触发。生产环境建议配置 \texttt{allowFrom} 白名单限制访问。
\end{quote}

\section{本章小结}

本章通过一个完整的实战案例，展示了在腾讯轻量云服务器上从零部署 OpenClaw 的全过程：

\begin{itemize}
  \item 从服务器环境准备、Node.js 安装到 OpenClaw 部署
  \item API 密钥配置、模型认证、守护进程设置
  \item 通过 TUI 交互模式体验 Agent 的完整能力
  \item 飞书渠道接入与配置
\end{itemize}

\noindent 至此，本书已经完整覆盖了 OpenClaw 的核心概念与实战操作。希望这份指南能帮助你快速上手 OpenClaw，
打造属于自己的个人 AI 智能体。
